% بانک تمرین‌ها — فصل 4
% فصل از نام این پرونده تعیین می‌شود.

\begin{exo}[انتگرال‌گیری از یک صورت دیفرانسیلی در امتداد مسیر]{integrale-forme-differentielle}
\keywords{صورت دیفرانسیلی، دیفرانسیل کامل، معیار شوارتس، انتگرال خطی}

دو صورت دیفرانسیلی زیر را در نظر می‌گیریم:
\[
\omega_1 = y\,\mathrm{d}x + x\,\mathrm{d}y,
\qquad
\omega_2 = y\,\mathrm{d}x,
\]
همچنین سه مسیر که مبدأ $O(0,0)$ را به نقطهٔ $M(1,1)$ وصل می‌کنند:
مسیر $\gamma_1$ از پاره‌خط افقی $O\to(1,0)$ و سپس پاره‌خط عمودی
$(1,0)\to M$ تشکیل شده است؛ مسیر $\gamma_2$ از پاره‌خط عمودی $O\to(0,1)$
و سپس پاره‌خط افقی $(0,1)\to M$؛ و مسیر $\gamma_3$ پاره‌خط مستقیمی است که
$O$ را مستقیماً به $M$ وصل می‌کند.

\begin{enumerate}
\item با یافتن تابعی مانند $f$ که $\omega_1=\mathrm{d}f$ باشد، نشان دهید $\omega_1$ کامل است.
\item با پارامتری‌کردن هر پاره‌خط، $\int_{\gamma_1}\omega_1$ و $\int_{\gamma_2}\omega_1$ را حساب کنید. نتیجه را بدون محاسبه نیز به دست آورید.
\item معیار شوارتس را بر $\omega_2$ اعمال کنید. چه نتیجه‌ای می‌توان گرفت؟
\item $\int_{\gamma_1}\omega_2$، $\int_{\gamma_2}\omega_2$ و $\int_{\gamma_3}\omega_2$ را حساب و تفسیر کنید.
\end{enumerate}

\begin{indication}
روی پاره‌خط افقی $\mathrm{d}y=0$ و روی پاره‌خط عمودی $\mathrm{d}x=0$ است؛
پس هر انتگرال به یک انتگرال ساده فروکاسته می‌شود.
\end{indication}

\begin{solution}
\textbf{پرسش 1.} تابع $f(x,y)=xy$ مناسب است، زیرا $\partial f/\partial x=y$ و $\partial f/\partial y=x$.

\textbf{پرسش 2.} برای منحنی پارامتری
$t\mapsto\bigl(x(t),y(t)\bigr)$ داریم
\[
\int_\gamma\omega_1
= \int \left[y(t)x'(t)+x(t)y'(t)\right]\,\mathrm{d}t.
\]

مسیر $\gamma_1$ اجتماع دو پاره‌خط است. روی نخستین، برای $t\in[0,1]$،
$x(t)=t$ و $y(t)=0$؛ پس $\mathrm{d}x=\mathrm{d}t$ و $\mathrm{d}y=0$.
روی دومی نیز برای $t\in[0,1]$، $x(t)=1$ و $y(t)=t$؛ پس
$\mathrm{d}x=0$ و $\mathrm{d}y=\mathrm{d}t$. بنابراین
\[
\begin{aligned}
I_1
= \int_{\gamma_1}\omega_1
&= \int_0^1\left(0\times 1+t\times 0\right)\,\mathrm{d}t
 + \int_0^1\left(t\times 0+1\times 1\right)\,\mathrm{d}t \\
&= 0+\int_0^1\mathrm{d}t = 1.
\end{aligned}
\]

برای $\gamma_2$، پاره‌خط نخست با $x(t)=0$ و $y(t)=t$ و دومی با
$x(t)=t$ و $y(t)=1$ پارامتری می‌شوند؛ در هر دو حالت $t\in[0,1]$. پس
\[
\begin{aligned}
I_2
= \int_{\gamma_2}\omega_1
&= \int_0^1\left(t\times 0+0\times 1\right)\,\mathrm{d}t
 + \int_0^1\left(1\times 1+t\times 0\right)\,\mathrm{d}t \\
&= 0+\int_0^1\mathrm{d}t = 1.
\end{aligned}
\]
دو مقدار برابرند. بدون محاسبه نیز می‌شد آن را پیش‌بینی کرد: چون
$\omega_1=\mathrm{d}f$ با $f(x,y)=xy$، انتگرال فقط به نقاط ابتدا و انتها بستگی دارد،
\[
\int_\gamma\omega_1=f(M)-f(O)=f(1,1)-f(0,0)=1.
\]

\textbf{پرسش 3.} می‌نویسیم $\omega_2=A(x,y)\,\mathrm{d}x+B(x,y)\,\mathrm{d}y$.
اینجا $A(x,y)=y$ و $B(x,y)=0$. اگر صورت کامل بود، معیار شوارتس می‌داد
\[
\frac{\partial A}{\partial y}=\frac{\partial B}{\partial x}.
\]
اما $\partial A/\partial y=1$، درحالی‌که $\partial B/\partial x=0$ است.
مشتق‌های متقاطع متفاوت‌اند؛ پس $\omega_2$ کامل نیست. انتگرال آن می‌تواند به
مسیر پیموده‌شده میان $O$ و $M$ وابسته باشد، چنان‌که محاسبهٔ بعدی نشان می‌دهد.

\textbf{پرسش 4.} چون $\omega_2=y\,\mathrm{d}x$، تنها تغییر $x$ در مقداری
غیرصفر از $y$ می‌تواند در انتگرال سهم داشته باشد. با پارامتری‌سازی‌های پیشین،
روی $\gamma_1$ داریم
\[
\begin{aligned}
J_1
=\int_{\gamma_1}\omega_2
&=\int_0^1 0\times 1\,\mathrm{d}t
  +\int_0^1 t\times 0\,\mathrm{d}t \\
&=0.
\end{aligned}
\]
روی پاره‌خط نخست $\gamma_2$، $x(t)=0$ و بنابراین $\mathrm{d}x=0$؛ روی
دومی، $x(t)=t$، $y(t)=1$ و $\mathrm{d}x=\mathrm{d}t$. ازاین‌رو
\[
\begin{aligned}
J_2
=\int_{\gamma_2}\omega_2
&=\int_0^1 t\times 0\,\mathrm{d}t
  +\int_0^1 1\times 1\,\mathrm{d}t \\
&=1.
\end{aligned}
\]
سرانجام پاره‌خط قطری $\gamma_3$ آشکارا چنین پارامتری می‌شود:
\[
x(t)=t,\qquad y(t)=t,\qquad t\in[0,1],
\]
پس $\mathrm{d}x=\mathrm{d}t$ و $\mathrm{d}y=\mathrm{d}t$. نتیجه می‌شود
\[
J_3
=\int_{\gamma_3}\omega_2
=\int_0^1 y(t)x'(t)\,\mathrm{d}t
=\int_0^1 t\,\mathrm{d}t
=\frac12.
\]
سه مسیر ابتدا و انتهای یکسان دارند، اما سه مقدار متفاوت می‌دهند:
$J_1=0$، $J_2=1$ و $J_3=1/2$. این دقیقاً وابستگی به مسیرِ مشخصهٔ یک صورت
دیفرانسیلی ناکامل است.
\end{solution}
\end{exo}

\begin{exo}[تغییر یکسان انرژی درونی، سه شیوهٔ انتقال]{meme-delta-u-trois-transferts}
\keywords{قانون اول، انرژی درونی، گرما، کار مکانیکی، کار الکتریکی، گرماسنجی، تابع حالت}

یک مایع، گرماسنج آن و مقاومت کوچکی که در آن غوطه‌ور است، سامانه‌ای بسته و
از دید ماکروسکوپی ساکن می‌سازند. ظرف صلب است و ظرفیت گرمایی کل سامانه ثابت
فرض می‌شود:
\[
C=2{,}00\ \mathrm{kJ\,K^{-1}}.
\]
می‌خواهیم سامانه را از حالت تعادل $A$ با $T_A=20{,}0\,^\circ\mathrm{C}$
به حالت تعادل $B$ با $T_B=25{,}0\,^\circ\mathrm{C}$ ببریم. می‌پذیریم که
\[
U(B)-U(A)=C(T_B-T_A).
\]
تغییرات انرژی جنبشی و پتانسیل ماکروسکوپی صفرند و همهٔ اتلاف‌ها به بیرون نادیده گرفته می‌شوند.

سه روش زیر جداگانه و از حالت اولیهٔ یکسان $A$ انجام می‌شوند:
\begin{itemize}
\item[روش 1] سامانه بدون دریافت هیچ کاری در تماس گرمایی با یک منبع گرم قرار می‌گیرد.
\item[روش 2] دیواره‌ها عایق گرمایی‌اند. یک پره به کمک سقوط جرم $m$ از ارتفاع $h=5{,}00$ m مایع را هم می‌زند. پره و مایع سرانجام ساکن می‌شوند و همهٔ انرژی مکانیکی ازدست‌رفتهٔ جرم به سامانه منتقل می‌شود. $g=9{,}81\ \mathrm{m\,s^{-2}}$ بگیرید.
\item[روش 3] سامانه گرمای $Q_3=4{,}00$ kJ دریافت می‌کند و باقی انرژی از راه الکتریکی به مقاومت داده می‌شود. توان الکتریکی دریافتی ثابت و برابر $\mathcal P=300$ W است.
\end{itemize}

\begin{enumerate}
\item تغییر انرژی درونی $\Delta U=U(B)-U(A)$ را که برای هر سه روش یکسان است حساب کنید.
\item برای هر یک از دو روش نخست، گرمای $Q$ و کار $W$ دریافتی سامانه را بیابید. در روش دوم جرم لازم $m$ را حساب کنید.
\item برای روش سوم، کار الکتریکی دریافتی و مدت برق‌رسانی به مقاومت را حساب کنید.
\item سه تحول حالت‌های ابتدا و انتهای یکسان دارند. توضیح دهید چرا بااین‌حال مقادیر $Q$ و $W$ می‌توانند متفاوت باشند. آیا سامانه در حالت نهایی $B$ «گرما» یا «کار» در خود دارد؟
\end{enumerate}

\begin{indication}
$\Delta U=Q+W$ را با قرارداد علامت بانکدار به کار ببرید. کار دریافتی پره برابر
کاهش $mgh$ در انرژی پتانسیل جرم است. انرژی الکتریکی که از راه سیم‌ها از مرز
می‌گذرد، کار الکتریکی شمرده می‌شود.
\end{indication}

\begin{solution}
\textbf{پرسش 1.} اختلاف دما $T_B-T_A=5{,}0$ K است. بنابراین
\[
\Delta U=C(T_B-T_A)
=2{,}00\times5{,}0
=10{,}0\ \mathrm{kJ}.
\]
این مقدار فقط به دو حالت تعادل $A$ و $B$ بستگی دارد.

\textbf{پرسش 2.} در روش نخست هیچ کاری دریافت نمی‌شود: $W_1=0$. قانون اول می‌دهد
\[
Q_1=\Delta U=10{,}0\ \mathrm{kJ}.
\]
گرما مثبت است، زیرا انرژی وارد سامانه می‌شود.

در روش دوم دیواره‌ها عایق گرمایی‌اند، پس $Q_2=0$. همهٔ تغییر انرژی درونی از
کار مکانیکی پره می‌آید:
\[
W_2=\Delta U=10{,}0\ \mathrm{kJ}.
\]
از $W_2=mgh$ به دست می‌آید
\[
m=\frac{W_2}{gh}
=\frac{10{,}0\times10^3}{9{,}81\times5{,}00}
\simeq 204\ \mathrm{kg}.
\]
بزرگی این مرتبه یادآور می‌شود که حتی افزایش اندک دما با انرژی مکانیکی
قابل‌توجهی متناظر است.

\textbf{پرسش 3.} قانون اول ایجاب می‌کند
\[
W_3=\Delta U-Q_3=10{,}0-4{,}00=6{,}00\ \mathrm{kJ}.
\]
این کار مثبت است: منبع برق به سامانه انرژی می‌دهد. با $W_3=\mathcal P\,\Delta t$،
\[
\Delta t=\frac{W_3}{\mathcal P}
=\frac{6{,}00\times10^3}{300}
=20{,}0\ \mathrm{s}.
\]

\textbf{پرسش 4.} سه مسیر به‌ترتیب می‌دهند
\[
(Q_1,W_1)=(10{,}0,0)\ \mathrm{kJ},
\qquad
(Q_2,W_2)=(0,10{,}0)\ \mathrm{kJ},
\]
و
\[
(Q_3,W_3)=(4{,}00,6{,}00)\ \mathrm{kJ}.
\]
در همهٔ حالت‌ها مجموع $Q+W$ برابر $10{,}0$ kJ است و همان تغییر $\Delta U$
را ایجاد می‌کند. انرژی درونی تابع حالت است، درحالی‌که گرما و کار انتقال‌های
انرژی در طول تحول را توصیف می‌کنند. در حالت نهایی $B$ سامانه انرژی درونی
دارد، اما «گرما» یا «کار» در خود ندارد.
\end{solution}
\end{exo}

\begin{exo}[تراکم زیر فشار خارجی ثابت]{compression-pression-exterieure-constante}
\seoready{true}
\keywords{کار نیروهای فشار، فشار خارجی ثابت، تراکم، انبساط، انبساط در خلأ}

گازی زیر فشار خارجی ثابت $P_0$ از حجم $V_A$ تا حجم $V_B<V_A$ متراکم می‌شود.

\begin{enumerate}
\item کار دریافتی گاز را حساب کنید.
\item گاز در برابر همان فشار خارجی $P_0$ از $V_B$ به $V_A$ بازمی‌گردد. کار دریافتی را حساب و با کار تراکم مقایسه کنید.
\item انبساط از $V_B$ به $V_A$ را هنگامی که گاز در خلأ منبسط می‌شود دوباره بررسی کنید. سه علامت به‌دست‌آمده را تفسیر کنید.
\end{enumerate}

\begin{indication}
از $W=-\int P_{\mathrm{ext}}\,\mathrm{d}V$ استفاده کنید. حتی وقتی تحول شبه‌ایستا
نیست، باید فشار \emph{خارجی} را انتگرال گرفت.
\end{indication}

\begin{solution}
\textbf{پرسش 1.} فشار خارجی ثابت است، پس
\[
W_{A\to B}=-P_0(V_B-V_A)=P_0(V_A-V_B)>0.
\]
گاز هنگام تراکم کار دریافت می‌کند.

\textbf{پرسش 2.} برای مسیر بازگشت،
\[
W_{B\to A}=-P_0(V_A-V_B)<0.
\]
دو کار مخالف یکدیگرند، زیرا دو تحول در برابر همان فشار خارجی ثابت انجام می‌شوند.

\textbf{پرسش 3.} در خلأ $P_{\mathrm{ext}}=0$ است، بنابراین
\[
W_{B\to A}^{\mathrm{vide}}=0.
\]
پس انبساط خودبه‌خود به معنی کار منفی نیست؛ باید واقعاً نیرویی خارجی به عقب رانده شود.
\end{solution}
\end{exo}
